\documentclass[12pt,a4paper]{article}

\usepackage[margin=2.5cm]{geometry}
\usepackage{booktabs}
\usepackage{tabularx}
\usepackage{array}
\usepackage{hyperref}
\usepackage{enumitem}
\usepackage{fancyhdr}
\usepackage{parskip}
\usepackage{titlesec}
\usepackage{xcolor}
\usepackage{longtable}

\definecolor{accent}{RGB}{30,80,160}

\hypersetup{colorlinks=true, linkcolor=accent, urlcolor=accent}

\titleformat{\section}{\large\bfseries\color{accent}}{}{0em}{}[\titlerule]
\titleformat{\subsection}{\normalsize\bfseries}{}{0em}{}

\pagestyle{fancy}
\fancyhf{}
\lhead{\textbf{WanderSync} \textit{Task 2: Stakeholder Identification}}
\rhead{Team 15}
\rfoot{\thepage}
\renewcommand{\headrulewidth}{0.4pt}

\newcolumntype{L}[1]{>{\raggedright\arraybackslash}p{#1}}

\begin{document}

\begin{center}
  {\LARGE\bfseries Task 2: Stakeholder Identification}\\[0.4em]
  {\large WanderSync: Dynamic Itinerary, Disruption Management and Social Travel Hub}\\[0.2em]
  {\normalsize Team 15}
\end{center}

\vspace{1em}

% ================================================================
\section{IEEE 42010 Architecture Framework}
% ================================================================

This document applies the IEEE/ISO/IEC 42010:2011 standard for systems and software engineering
architecture descriptions. The standard defines a systematic way to document the architecture of a
system by identifying \textit{stakeholders}, their \textit{concerns}, the \textit{viewpoints} that
frame those concerns, and the \textit{views} that instantiate each viewpoint for WanderSync.

% ================================================================
\section{Stakeholders and Concerns}
% ================================================================

An architecture description must address the concerns of all relevant stakeholders. The following
table enumerates the stakeholders of WanderSync and the architectural concerns each holds. A concern
is any interest that a stakeholder has in the system that is architecturally relevant.

\vspace{0.5em}
\begin{tabularx}{\textwidth}{L{3.8cm}X}
\toprule
\textbf{Stakeholder} & \textbf{Architectural Concerns} \\
\midrule
\textbf{Traveller} (end user) &
  Accuracy and completeness of the unified itinerary timeline; timely delivery of disruption
  notifications; enforcement of location privacy before bilateral consent; learnability and
  responsiveness of the interface (3-second page load). \\
\addlinespace
\textbf{Development Team} &
  Modifiability and low coupling (new booking provider must touch $\leq$2 files outside the
  adapter module); testability of core service functions ($\geq$60\% line coverage); clear module
  boundaries for parallel development; onboarding speed (new team member can run the system
  locally within 15 minutes). \\
\addlinespace
\textbf{System Administrator / DevOps} &
  Deployability on standard hosting infrastructure; availability and crash-free operation during
  all planned demonstration scenarios; observability through structured logging; straightforward
  configuration of environment variables and secrets. \\
\addlinespace
\textbf{Travel Provider} (external API) &
  Stability of the integration contract exposed by the adapter interface; isolation from changes
  in other providers; well-defined error propagation so that a provider outage does not cascade. \\
\addlinespace
\textbf{Privacy Regulator / Compliance} &
  Data minimisation: no precise location stored beyond the active session unless explicitly saved;
  zero location leakage before bilateral opt-in; consent state enforced at the data-access layer,
  not only the UI; all credentials handled with bcrypt and transmitted over HTTPS. \\
\bottomrule
\end{tabularx}

% ================================================================
\section{Viewpoints}
% ================================================================

A \textit{viewpoint} is a specification of the conventions used to construct and use a view. Each
viewpoint below names the stakeholder concerns it addresses and describes the modelling conventions
and notation used in its corresponding view.

\subsection{Viewpoint 1: Functional Viewpoint}

\textbf{Addresses concerns of:} Traveller, Development Team.

\textbf{Purpose:} Describes the functional decomposition of WanderSync into subsystems and the
responsibilities and interfaces of each. Answers the question: \textit{what does the system do and
how is that responsibility partitioned?}

\textbf{Modelling conventions:} UML component diagram showing subsystems as components, provided
and required interfaces, and connector types (REST, message broker, direct function call).

\subsection{Viewpoint 2: Information Viewpoint}

\textbf{Addresses concerns of:} Traveller, Privacy Regulator.

\textbf{Purpose:} Describes the key data entities managed by the system, their ownership, lifecycle,
and the rules governing their disclosure. Answers: \textit{what data does the system hold, and how
does it flow and transform?}

\textbf{Modelling conventions:} Entity–relationship overview of \texttt{BookingRecord},
\texttt{Itinerary}, \texttt{DisruptionEvent}, \texttt{User}, and \texttt{ConsentState}; data-flow
arrows annotated with privacy rules (e.g., city-level only until consent $=$ \textit{precise}).

\subsection{Viewpoint 3: Deployment Viewpoint}

\textbf{Addresses concerns of:} System Administrator, Development Team.

\textbf{Purpose:} Describes the runtime topology — how system components are assigned to execution
environments and how they communicate across network boundaries. Answers: \textit{where does the
system run and how are processes connected?}

\textbf{Modelling conventions:} UML deployment diagram showing nodes (client browser, API Gateway,
backend service processes, databases, message broker) and communication paths with protocol
annotations.

\subsection{Viewpoint 4: Security Viewpoint}

\textbf{Addresses concerns of:} Privacy Regulator, Traveller, System Administrator.

\textbf{Purpose:} Describes how authentication, authorisation, data-in-transit protection, and
location-disclosure control are enforced across the system. Answers: \textit{how does the system
protect user identity and sensitive data?}

\textbf{Modelling conventions:} UML sequence diagram of the JWT authentication flow; state diagram
of the consent state machine (\textit{none} $\to$ \textit{city-level} $\to$ \textit{precise});
annotations of TLS termination points.

% ================================================================
\section{Views}
% ================================================================

Each view below instantiates its corresponding viewpoint for the WanderSync system. Where a
diagram is specified, the PlantUML source is maintained in the \texttt{diagrams/} directory of
this deliverable and must be rendered to PNG before final submission.

\subsection{View 1: Functional View — Subsystem Component Diagram}

WanderSync is decomposed into six top-level subsystems connected through explicit interfaces:

\begin{enumerate}[noitemsep]
  \item \textbf{API Gateway \& Frontend Layer} — single external entry point; TLS termination;
        JWT validation middleware; React SPA delivery.
  \item \textbf{Authentication \& User Management Service} — registration, login, JWT issuance,
        profile CRUD.
  \item \textbf{Booking Aggregation Service} — \texttt{BookingAdapter} per provider; normalises
        heterogeneous provider data to \texttt{BookingRecord}.
  \item \textbf{Itinerary Management Service} — authoritative itinerary store; timeline assembly;
        caching; consumes \texttt{DisruptionEvent} messages.
  \item \textbf{Disruption Detection Service} — monitors external feeds; publishes
        \texttt{DisruptionEvent} to message broker.
  \item \textbf{Social Synchronisation Service} — geospatial matching; bilateral consent state
        machine; enforces privacy-preserving disclosure rules.
\end{enumerate}

The Disruption Detection Service and Itinerary/Notification services communicate exclusively
through the message broker, realising the event-driven decoupling documented in ADR-001.

\subsection{View 2: Information View — Data Entities and Privacy Rules}

The core data entities and their disclosure constraints are:

\begin{tabularx}{\textwidth}{L{3.5cm}L{4cm}X}
\toprule
\textbf{Entity} & \textbf{Owner} & \textbf{Privacy / Lifecycle Rule} \\
\midrule
\texttt{User} & Auth Service & Passwords stored as bcrypt hash; never exposed raw. \\
\texttt{BookingRecord} & Booking Aggregation & Normalised from provider data; visible only to
  owning user. \\
\texttt{Itinerary} & Itinerary Service & Assembled from BookingRecords; cached; updated on
  DisruptionEvent. \\
\texttt{DisruptionEvent} & Disruption Service & Published to broker; consumed by Itinerary and
  Notification services; not persisted long-term. \\
\texttt{ConsentState} & Social Service & Values: \textit{none}, \textit{city-level},
  \textit{precise}; transitions require explicit bilateral action; governs location resolution
  returned by the Social Service. \\
\texttt{LocationData} & Social Service & Raw coordinates never persisted beyond session unless
  user explicitly saves; only city-level label exposed when ConsentState $\neq$ \textit{precise}. \\
\bottomrule
\end{tabularx}

\subsection{View 3: Deployment View — Runtime Topology}

The runtime topology comprises the following nodes and communication paths:

\begin{itemize}[noitemsep]
  \item \textbf{Client Browser} — communicates with API Gateway over HTTPS (port 443).
  \item \textbf{API Gateway} — proxies to Auth, Itinerary, Disruption, and Social services over
        internal HTTP; enforces JWT before forwarding.
  \item \textbf{Backend Services} — Auth, Booking Aggregation, Itinerary Management, Disruption
        Detection, Social Synchronisation — each runs as an independent process.
  \item \textbf{Message Broker} — Disruption Detection publishes; Itinerary and Notification
        services subscribe.
  \item \textbf{Databases} — User DB (Auth Service); Itinerary DB (Itinerary Service); Social DB
        (Social Service). Each service owns its own store.
  \item \textbf{External APIs} — Flight tracking and weather services polled or subscribed to by
        the Disruption Detection Service over HTTPS.
\end{itemize}

\subsection{View 4: Security View — Authentication and Consent}

\textbf{JWT Authentication Flow:}
\begin{enumerate}[noitemsep]
  \item Client submits credentials to API Gateway $\to$ forwarded to Auth Service.
  \item Auth Service verifies bcrypt hash; issues signed JWT (HS256, 24-hour expiry).
  \item JWT returned to client; stored in memory (not localStorage).
  \item All subsequent requests carry JWT in \texttt{Authorization: Bearer} header.
  \item API Gateway validates JWT signature and expiry before proxying; rejects with 401 on
        failure.
\end{enumerate}

\textbf{Location Consent State Machine:}
\begin{itemize}[noitemsep]
  \item Initial state: \textit{none} — only city-level label returned; no coordinates.
  \item User A sends opt-in request $\to$ state transitions to \textit{pending} on A's side.
  \item User B accepts $\to$ bilateral consent established; state becomes \textit{precise} for
        both.
  \item Either user revokes $\to$ state resets to \textit{none}; precise coordinates withheld
        immediately.
  \item Enforcement is applied at the data-access layer of the Social Service, not the UI.
\end{itemize}

\end{document}
